% !TEX root = main.tex

For many decades, graph theory has been a tool used to model and study many real
world problems and phenomena~\cite{BondyM76}. A usual assumption for many of
these problems is that the graphs have a fixed structure. However, there are
quite a number of cases where the structure of a system changes over time. For
example, consider the problem of routing in transportation networks (roads,
rails) where specific connections can become unavailable (due to a disaster) or
they are active during specific times (due to safety). Such scenarios can be
modeled using temporal graphs, a sequence of graphs over the same vertex set
where the edges possibly change in each time step. The temporal graph setting
has received significant interest from the research community in the recent past
as seen in recent surveys~\cite{CasteigtsFQS12,Michail16}. 

In this paper, we study two problems on temporal graphs. The first is the
\textit{temporal exploration problem} ($\texp$), which has been studied, e.g., by Michail and
Spirakis~\cite{MichailS14,MichailS16} and by Ilcinkas et
al.~\cite{IlcinkasKW14}. It requires an agent to explore all vertices of the
temporal graph as quickly as possible. The second is the \textit{temporal
rendezvous problem} ($\trp$), which we formulate for the first time in this
paper. It requires two heterogeneous agents (in terms of the programs they run)
to rendezvous on a temporal graph when they cannot communicate with one another.
For both problems, we assume that each  agent has complete knowledge of the
temporal graph in advance (a common
assumption~\cite{AdamsonGMZ22,BodlaenderV19,ErlebachHK21,ErlebachKLSS19,ErlebachS22a,ErlebachS22b,IlcinkasKW14,MichailS14,MichailS16,Taghian20}).
However, in the case of $\trp$ the agents may have different names for the same
vertices, i.e., the local labels of the vertices may be different. 

The problem of exploration has been well studied in the static setting since it was introduced in 1951 by Shannon~\cite{Shannon51}.  
It has also been intensively studied in the temporal graph setting since
2014 (see all references from the previous paragraph).
%~\cite{AdamsonGMZ22,BodlaenderV19,ErlebachHK21,ErlebachKLSS19,ErlebachS22a,ErlebachS22b,IlcinkasKW14,MichailS14,MichailS16,Taghian20}.
On an application-oriented note, the problem captures the setting where a person
is trying to visit various parts of a city using %constantly 
public
transportation. For example, train schedules involve multiple train
stations (vertices) with trains running between them at different times (i.e.,
%constantly 
repeatedly
changing edges). Thus, planning a visit to multiple destinations
over a given day using railways is an example of solving $\texp$.


The rendezvous problem can be broadly categorized into two types: symmetric and
asymmetric rendezvous. The version where agents have the same strategy
(symmetric rendezvous) was introduced by Alpern~\cite{Alpern76}. The version
where agents can have distinct strategies (asymmetric rendezvous) was introduced
by Alpern~\cite{Alpern95} and is the focus of research in this paper. 
%Both
%versions have received much interest over the years resulting in numerous
%surveys~\cite{Alpern02,AlpernG06,Flocchini19,KranakisKM22,KranakisKR06,Pelc12,Pelc19}
%being written on the problems covering different settings.  
$\trp$ is a natural
extension of the asymmetric rendezvous problem to the dynamic setting. As a real
world example, consider a pair of tourists who want to explore a city together
and have to agree on a strategy to meet up in case they are separated and their
cell phones die. In this scenario, they may use public transportation
(dynamically changing network) to meet and agree in advance to use different
strategies that guarantee that they meet quickly. 

In this paper, we present results that extend the literature in two ways.
Firstly, we formalize the \trp{} problem in the setting where agents have
complete knowledge of the temporal graph a priori and develop good upper and
lower bounds for it. 
Secondly, we
utilize interesting structural properties (namely automorphisms of a graph and
associated orbits of  vertices) in order to analyze 
bounds on the number of time steps required by temporal walks that solve certain problems.
In particular, we develop upper and lower bounds for both \texp{} and \trp{} that leverage the 
aforementioned graph properties. To the best of our knowledge, this is the first work that takes 
advantage of these graph properties to study problems in temporal graphs. 

\subsection{Our Contributions}
\label{subsec:our-contribs}
We present results for two problems: first, the \textsc{temporal exploration
problem} ($\texp$) which, for a given temporal graph $\mathcal{G}$, asks for a
temporal walk that visits all vertices of $\mathcal{G}$, and secondly, the
\textsc{temporal rendezvous problem} (\trp), which considers two agents that try
to meet in the given temporal graph, meaning they must be stationed at the same
vertex in the same time step.

One of our primary contributions is formalizing the problem of $\trp$ in
Section~\ref{sec:pre}, and in doing so extending the problem of asymmetric
rendezvous to the temporal graph setting where agents have complete knowledge of
the temporal graph in advance. {Another significant contribution is that we
show} how to leverage the use of a structural graph property, namely the
automorphism group of a temporal
graph and the associated notion of orbits, to bound the number of time steps
required by algorithms we devise for these problems. To the best of our
knowledge, this work is the first instance of leveraging such properties to
study any problem on temporal graphs apart from a practical
implementation of a generator for listing all non-isosomorphic simple temporal
graphs~\cite{Casteigts20}. Intuitively, an automorphism of a graph is a mapping
from the set of vertices of the graph to the same set of vertices that preserves
the
%notion of ``neighborliness''
neighborhood relation between the vertices. The set of automorphisms of a temporal graph
consists of the \textit{intersection} of the sets of automorphisms of the
graph at each time step. The set of automorphisms of the temporal graph,
along with function composition as group operation, forms the automorphism group of the
temporal graph. An orbit of the automorphism group of a temporal graph is a maximal set
of vertices such that each vertex can be mapped to any other vertex in the set via
an automorphism of the group.
Intuitively, the vertices in an orbit look indistinguishable
to an agent that has full information about the temporal graph but
without meaningful vertex labels. As the agents are able to compute a
consistent numbering of the orbits (Lemma~\ref{lem:meeting}), they can agree to meet
in some specific orbit, but not at a specific vertex.
{Therefore, temporal graphs where all orbits are large
(or that even have a single orbit containing all vertices)
appear to be the most challenging graphs for solving
$\trp$. Our result providing fast exploration schedules
in temporal graphs with few orbits is therefore a
crucial ingredient for enabling our solution to $\trp$ to handle
all possible temporal graphs (including those with a single
orbit).}


%
%  already said above: "we formally define $\trp$ in Section~\ref{sec:pre}"
%
%We introduce several useful definitions and formally define
%$\texp$ and $\trp$ in Section~\ref{sec:pre}. 
%

We give precise definitions and present further preliminaries in
Section~\ref{sec:pre}. Then,
we introduce some useful
utilities related to automorphisms in Section~\ref{sec:aut-util}. These
will be used in later sections, where we present the following
results. 

\begin{sloppypar}
\textbf{Upper bounds.} 
In Section~\ref{sec:texp}, we develop a deterministic algorithm to solve $\texp$
%in $O(\min\{n^2,r n^{1+\epsilon}\})$ time steps for any fixed $\epsilon >0$ (see
in $O(r n^{1+\epsilon})$ time steps for any fixed $\epsilon >0$ (see
Corollary~\ref{cor:epsilonallorbits}), where $n$ is the number of vertices in
the temporal graph and $r$ is the number of orbits of the automorphism group of
the temporal graph. Note that $r$ can range in value from $1$ to $n$. Thus, for
$r=O(n^c)$ for constant $c<1$, this is a significant improvement over the known
tight worst-case bound of $\Theta(n^2)$ time steps for
arbitrary connected temporal graphs~\cite{ErlebachHK21}. %Subsequently 
In Section~\ref{sec:rendezvous}, {we leverage this algorithm for $\texp$ to
develop a deterministic solution for $\trp$ using} $O(n^{1+\epsilon})$ time
steps for any fixed $\epsilon > 0$ (see Theorem~\ref{thm:rendezvous-time}). 
Our focus is on bounding the time steps of the temporal walks required to
solve $\texp$ and $\trp$, and not on \mbox{optimizing the running time of the
respective algorithms to compute such walks.}
\end{sloppypar}

\textbf{Lower bounds.} 
We complement our algorithms with lower bounds for both $\texp$ and $\trp$ in
Section~\ref{sec:lower-bounds}. In particular, we design an instance of $\trp$
such that any solution for it requires $\Omega(n \log n)$ time steps (see
Theorem~\ref{thm:rendezvouslb}). We then show how this translates to a lower
bound of $\Omega(n \log n)$ time steps for {some instances of $\texp$ where
the temporal graph at hand has an orbit number $r=O(1)$ (see
Corollary~\ref{cor:explorationlb})}. By revisiting the lower bound of
\cite{ErlebachHK21} for $\texp$, which focused on arbitrary temporal graphs, and
studying it through the lens of automorphisms and orbits, we can obtain a more
fine-grained lower bound of $\Omega(rn)$ time steps (see
Lemma~\ref{lem:rn-lower-bound}) for {some temporal graphs with orbit number
$1\le r \le n$.} Notice that the {multiplicative} gap between our upper and lower bounds for both
problems is only {a factor of} $O(n^\epsilon)$.

{ \textbf{Relevance of the orbit number.} {It is well
known in graph theory}
    that almost all (static) graphs are \textit{rigid}~\cite{BallGS18},
    meaning that the automorphism group contains only the trivial
    identity function. This is in contrast to observations in practice.
    Many real-world graphs have non-trivial automorphisms. A recent
    analysis of real-world graphs in the popular database
    \texttt{networkrepository.com} showed that over $70\%$ of the analyzed
    graphs had non-trivial automorphisms~\cite{BallGS18}. One may
    reasonably expect that real-world temporal graphs have similar
    properties. Symmetries are also abundant in graphs arising in
    chemistry~\cite{Balasubramanian82}, which have been studied in temporal
    settings as well~\cite{PlamperLH23}.
	{We believe that these observations provide a strong
	motivation for studying temporal graph problems for temporal
	graphs with fewer than $n$ orbits,{\ e.g., with
    algorithms parameterized via the number of orbits.}}
    %
	{Furthermore, we emphasize that
    our results for $\trp$ {hold for all temporal graphs,}
	independent of the orbit number parameter,
even {though} they {utilize} techniques we develop for $\texp$
parameterized by the orbit number.} {Roughly speaking,
for orbit number~$r$, the penalty factor $r$ in the number of time steps 
to solve $\texp$ is saved when solving $\trp$ by focusing only on a smallest orbit of size at most $n/r$.
}
}

\subsection{Technical Overview and Challenges}
\label{subsec:tech-overview-challenges}


Firstly, we give an intuitive overview of our upper bound results. The concepts
used in this section are more precisely defined in the next sections.

For $\texp$, we first consider the problem of visiting all the vertices of one
orbit~$S$. One key insight is that, if we have a temporal walk $W$ that visits
$k$ vertices of~$S$, we can use the automorphisms of the temporal graph to
transform $W$ into other walks that visit different sets of $k$ vertices of~$S$
(Lemma~\ref{lem:transform1}).
Therefore, even if the $k$ vertices visited by $W$ have already been explored
earlier, we can transform $W$ into a temporal walk $W'$ that visits a ``good''
number of previously unexplored vertices of~$S$. The number of previously
unexplored vertices of $S$ that $W'$ is guaranteed to visit increases with the
number of possible start vertices in $S$ that we allow for $W'$, but the larger
that set $X$ of possible start vertices is, the longer it may take the agent to
reach the best start vertex in that set. A challenge is to analyze this
tradeoff. By carefully relating the size of $X$ to the guaranteed number of
unexplored vertices that a walk $W'$ starting at a vertex in~$X$ can visit
(Corollary~\ref{cor:transform2}), we
manage to balance the number of time steps needed to move to the start vertex of
$W'$ and the number of previously unexplored vertices that $W'$ visits.
To show that vertices of $X$ can be reached quickly, we study the structure
of edges that connect vertices in different orbits (Lemma~\ref{lem:orbitedges})
and use it to analyze
reachability between orbits (Lemmas~\ref{lem:laneorbit} and~\ref{lem:nextreachable}).

We then employ a recursive construction: We recursively construct a temporal
walk {$W_{1}$ that} visits $k$ vertices of~$S$, concatenate
$W_{1}$ with a temporal walk that moves quickly from the endpoint of
$W_{1}$ to a good start vertex in $S$ for what follows, and then use a second recursively constructed walk
(transformed via an automorphism to a ``good'' temporal walk
{$W_{2}$} starting in a vertex of $S$) to visit ``nearly'' $k$
further vertices of~$S$. {A careful} analysis then shows that in this way
we can visit a constant fraction of the vertices of $S$ in
$O(r|S|^{1+\epsilon})$ steps (Lemma~\ref{lem:vtexplorationub} and Corollary~\ref{cor:tinyfraction}).
We then show that the concatenation of $O(\log
|S|)$ such walks (each one again transformed via an automorphism to maximize the
number of newly explored vertices) suffices to visit all vertices of $S$ in
$O(r|S|^{1+\epsilon}+n\log |S|)$ time steps (Theorem~\ref{thm:epsilonexploration}),
which can be bounded by $O(r
n^{1+\epsilon'})$ time steps). By visiting the $r$ orbits one after another, we
can finally show that the whole temporal graph can be explored in
$O(r n^{1+\epsilon'})$ time steps (Corollary~\ref{cor:epsilonallorbits}).

For $\trp$, a simple and fast solution one first thinks of is to have the agents
simply meet at a vertex with a specific label. {However}, this is not
feasible in the model we consider. In particular, we assume that the agents
cannot communicate %\footnote{{While it is unclear whether the ability to communicate
%would enable asymptotically faster solutions to $\trp$, it could help the
%agents to meet faster in certain instances: For example, the two agents
%could inform each other about the orbit in which they are located
%{\color{blue}(we show in Lemma~\ref{lem:meeting} that the
%agents can identify orbits using consistent integer colorings)}, and
%if they happen to be located in the same orbit and that orbit is small,
%they could meet there
%instead of both moving to a {designated} orbit to meet.}{\ In detail, this saves
%the time steps required to move to the designated orbit.}}
and, while both agents have complete information about the
temporal graph, they do not have access to consistent vertex labels. As such,
the agents are unable to agree upon the same vertex based on the vertex label.
We rely, instead, on structural graph properties and let the agents meet at a
vertex in a smallest orbit: One agent moves to any vertex in that orbit, and the
other searches all vertices in that orbit. The first challenge is for the agents
to independently identify the same orbit in which to meet. We show that this can
be done by letting each agent enumerate all temporal graphs, together with
colorings of their orbits, until it encounters for the first time a temporal
graph that is isomorphic to the input graph in which $\trp$ is to be solved. In
this way both agents obtain the same colored temporal graph and can
independently select, among all orbits of smallest size, the one with smallest
color (Lemma~\ref{lem:meeting}).
%
Then one agent moves to a vertex in that orbit~$S$, while the other searches all
vertices of~$S$. The challenge here is to deal with orbits that are large,
because previously known techniques would require $\Theta(n^2)$ time steps to
explore an orbit of size $\Omega(n)$. {Fortunately}, this is where our
above-mentioned results for exploration of (orbits of) temporal graphs come to
the rescue: As $|S|\le n/r$, we have $O(r|S|^{1+\epsilon}+n\log
|S|)=O(n^{1+\epsilon'})$, and hence $S$ can be explored (and $\trp$ solved) in
$O(n^{1+\epsilon'})$ time steps (Theorem~\ref{thm:rendezvous-time}).

Now we turn to our lower bounds.
For $\texp$, we first observe that the
existing lower bound construction that shows that $\Omega(n^2)$ time steps are
necessary for exploration on some temporal graphs uses temporal graphs with
$\frac{n}{2}+1$ orbits. By slightly varying the construction, we can show
for any $r$ in the range from $1$ to $n$
that there are temporal graphs with $r$ orbits that require
$\Omega(rn)$ time steps for exploration (Lemma~\ref{lem:rn-lower-bound}).
Our main lower bound result is the lower bound of
$\Omega(n\log n)$ for $\trp$ in temporal graphs with a single orbit
(Theorem~\ref{thm:rendezvouslb}). The
temporal graph is a cycle {on $n=2^m-1$ vertices} in every step, and the
edges change every $\lfloor n/16\rfloor$ time steps. Each period of $\lfloor
n/16\rfloor$ time steps in which the edges do not change is called a
\emph{phase}. We number the vertices of the cycle from $0$ to $n-1$ and consider
the binary representation of the vertex labels. {In phase~$1$, each vertex
$j$ is adjacent to $j+1$ and $j-1$, while in any phase $i>1$, it is adjacent to
$j+2^{2i}$ and $j-2^{2i}$ (with all computations done modulo~$n$).} As all
vertices are indistinguishable to the agents, we can force each agent to make
the same movements no matter where we place it initially. By fixing the five
highest-order bits of the agents' start positions in a certain way, we can
ensure that the agents do not meet in the first phase, no matter how the
remaining bits of their start positions are chosen. Depending on the positions
where the agents end up at the end of each phase (relative to their start
position in that phase), {we fix a certain number of bits of the binary
representations of the start positions of the agents: After the first phase, we
fix the lowest four bits, and after each further phase, we fix the next-higher
two bits. We can show that this is sufficient to ensure} that the agents start
the next phase at vertices that are sufficiently far from each other in the
cycle of that phase. This process can be repeated $\Omega(\log n)$ times,
showing that it takes $\Omega(n \log n)$ {time steps} for the agents to meet. This lower
bound for $\trp$ implies also that exploration of the constructed temporal graph
requires $\Omega(n \log n)$ steps (Corollary~\ref{cor:explorationlb}).

\full{In Section~\ref{sec:rand} we present a randomized algorithm for
    computing a temporal walk $W$ that visits all vertices of a given
    orbit with probability at least $1-\epsilon$ for any $0<\epsilon<1$.
    The computed temporal walk $W$ spans $O((n^{5/3}+rn)\log n)$ time
    steps, with $r$ the number of orbits of the automorphism group. The
    walk can be computed by $O(n^{1/3} \log (n/\epsilon))$ linear-time
    scans of the snapshots that exist in the aforementioned time span. We
    also give an alternative algorithm that has an expected runtime that
    consists of $O(n^{1/3})$ linear-time scans, but provides the specified
    temporal walk with probability $1$. In contrast to the upper bounds
    presented in the preceding sections, the algorithm has a runtime that
is independent of the number of automorphisms of the temporal graph at
hand.}

\subsection{Related Work}
\label{subsec:related-work}

There is a divide in the research on temporal graphs with respect to the
amount of knowledge known in advance by the agents on the graph.  When
complete knowledge of the temporal graph is known in advance to the agents,
as is assumed in this paper, the setting is
sometimes referred to as the \emph{post-mortem}
setting (see, e.g., Santoro~\cite{Santoro15}). Since we consider
scenarios where agents plan temporal walks using advance knowledge
of future time steps, we refer to the post-mortem
setting as the {\em clairvoyant} setting in the remainder of this section.

The clairvoyant setting is in contrast to the \emph{live} setting where agents only have partial
knowledge of the temporal graph, and the solutions to problems in these
different settings are of a different nature.  The reader can refer to the
survey by Di Luna~\cite{DiLuna19} for more information on problems and
solutions in the live setting.  We restrict the rest of this related work
section to work in the clairvoyant setting.


~\\ \noindent \textbf{Exploration.} 
%
% known from the introduction
%
%The problem of exploration was introduced in
%1951 by Shannon~\cite{Shannon51} and has been well studied in various settings.
%The temporal graph exploration problem ($\texp$) has been studied intensively
%since 2014, e.g., by Michail and Spirakis~\cite{MichailS14,MichailS16} and Ilcinkas et
%al.~\cite{IlcinkasKW14} and has been studied under various
%aspects~\cite{AdamsonGMZ22,BodlaenderV19,ErlebachHK21,ErlebachKLSS19,ErlebachS22a,ErlebachS22b,Taghian20}.
We trace the progress of solving instances of $\texp$, first by identifying
which assumptions are required for the problem to even be solvable, and then by
identifying properties that were leveraged to give faster and faster solutions. 

Michail and Spirakis~\cite{MichailS14,MichailS16} showed that it is
\textbf{NP}-complete to decide whether a given temporal graph with a given
\textit{lifetime}, i.e., the number of time steps it exists, is explorable when
no assumptions are made on the input graph. This holds even if the graph is connected in
every time step, termed \textit{connected} (and sometimes called \emph{always-connected}
in the literature).
Even when a restriction is placed on the \textit{underlying graph}, i.e., the
union of the graphs at each time step, such that the underlying graph has
pathwidth at most~$2$, Bodlaender and van der Zanden~\cite{BodlaenderV19} showed
that $\texp$ is \textbf{NP}-complete. 

However, the exploration is always possible when the lifetime of the graph is
sufficiently large. In particular, Michail and Spirakis~\cite{MichailS16} showed
that a connected graph with $n$ vertices may be explored in $O(n^2)$ time steps. For
the rest of the related work on $\texp$, we focus on the case of connected
temporal graphs with a sufficiently large lifetime. Erlebach et
al.~\cite{ErlebachHK21} showed that exploration on arbitrary temporal graphs
takes $\Omega(n^2)$ time steps. By restricting their study of temporal
graphs to those where the underlying graph belongs to a special graph class,
however, they showed that $\texp$ can be solved in $o(n^2)$ time steps in
several such cases. In particular,
when the underlying graph is planar, has bounded treewidth~$k$, or is a $2
\times n$ grid, they showed that $\texp$ can be solved in $O(n^{1.8} \log n)$
time steps, $O(n^{1.5} k^{1.5} \log n)$ time steps, and $O(n \log^3 n)$ time
steps, respectively. They also showed a lower bound of $\Omega(n \log n)$ when
the underlying graph is a planar graph of degree at most $4$. 

The study of $\texp$ when the temporal graph is restricted continued in
several papers. Taghian Alamouti~\cite{Taghian20} showed that $\texp$ can be
solved in $O(k^2 (k!) (2e)^k n)$ time steps when the underlying graph is a cycle
with $k$ chords. Adamson et al.~\cite{AdamsonGMZ22} improved this to $O(kn)$
time steps. They also improved the upper bounds on $\texp$ for underlying graphs
that have bounded treewidth $k$ or are planar to $O(k n^{1.5} \log n)$ and
$O(n^{1.75} \log n)$, respectively. In addition, they strengthened the lower bound for
underlying planar graphs by showing that even if the degree is at most~$3$, the
lower bound is $\Omega(n \log n)$. Erlebach et al.~\cite{ErlebachKLSS19} further
improved work on bounded degree underlying graphs by showing that  $\texp$ can
be solved in $O(n^{1.75})$ time steps for such temporal graphs.
Ilcinkas et al.~\cite{IlcinkasKW14}
showed that when the underlying graph is a cactus, the exploration time is
$2^{\Theta(\sqrt{\log n})} n$ time steps. 
%, and show a lower bound of
%$2^{\Omega(\sqrt{\log n})} n$ time steps. 

Other variants of $\texp$ have been studied where the problem is slightly
different, or the edges of the temporal graph vary in some particular way (e.g.,
periodically, $T$-interval connected, $k$-edge deficient), or multiple agents
explore the
graph~\cite{AaronKM14a,AaronKM14b,AkridaMSR21,BumpusM23,ErlebachS22a,ErlebachS22b,IlcinkasW13}. 


~\\ \noindent \textbf{Rendezvous.} 
%The rendezvous problem has been studied
%extensively in static graphs and metric spaces, after symmetric rendezvous was
%introduced by Alpern~\cite{Alpern76} under a different name, and both it and
%asymmetric rendezvous were formalized by Alpern~\cite{Alpern95}. 
%
Symmetric rendezvous~\cite{Alpern76} and asymmetric rendezvous ~\cite{Alpern95}
have received much interest over the years, resulting in numerous
surveys~\cite{Alpern02,AlpernG06,Flocchini19,KranakisKM22,KranakisKR06,Pelc12,Pelc19}
being written on the problems covering different settings. 
In this work, we are
the first to extend asymmetric rendezvous to temporal graphs in the clairvoyant
setting. Note that there has been previous
work~\cite{BournatDP18,DasDPP19,DiLunaFPPSV20,MichailST21,OoshitaD18,ShibataKESNK23,ShibataSNK21}
that has studied multi-agent rendezvous (also called gathering) in a dynamic
graph setting, but that was in the live setting. 
%
For an overview of rendezvous in static graphs we refer the reader to~\cite{AlpernG06,
KranakisKR06}. Rendezvous has been studied in deterministic
settings~\cite{BhagatP22,Pelc23} and asynchronous
settings~\cite{dieudonne2013meet,MARCO2006315}, and has been shown to be
solvable in logarithmic space~\cite{CzyzowiczJKAP10}.

~\\ \noindent \textbf{Other Related Work.}  
Other problems have also been studied in the temporal graph setting, e.g.,
matchings~\cite{MertziosMNZZ20},
separators~\cite{FLUSCHNIK2020197,ZSCHOCHE202072}, vertex
covers~\cite{AKRIDA2020108,HammKMS22}, containments of
epidemics~\cite{ENRIGHT202160}, Eulerian tours~\cite{BumpusM23,MarinoS21}, graph
coloring~\cite{MarinoS22,MertziosMZ21}, network flow~\cite{AKRIDA201946},
treewidth~\cite{Fluschnik2020} and cops and robbers~\cite{Erlebach-Spooner/20,MorawietzRW20}. See
the survey by Michail~\cite{Michail16} for more information.
